% SMPdesign/partexercises.tex
% mainfile: ../perfbook.tex
% SPDX-License-Identifier: CC-BY-SA-3.0

\section{分区练习}
\label{sec:SMPdesign:Partitioning Exercises}
%
\epigraph{当一个理论在你看来是唯一可能的理论时，
	  这表明你既没有理解这个理论，
	  也没有理解它试图解决的问题。}
	  {Karl Popper}

尽管分区技术比2000年代早期更广为人知，但其价值仍然被低估。
因此，\Cref{sec:SMPdesign:Dining Philosophers Problem}
以更高度并行的视角审视经典的哲学家就餐问题，
而\cref{sec:SMPdesign:Double-Ended Queue}
则重新审视双端队列。

\subsection{哲学家就餐问题}
\label{sec:SMPdesign:Dining Philosophers Problem}

\begin{figure}
\centering
\includegraphics[scale=.7]{SMPdesign/DiningPhilosopher5}
\caption{哲学家就餐问题}
\ContributedBy{fig:SMPdesign:Dining Philosophers Problem}{Kornilios Kourtis}
\end{figure}

\Cref{fig:SMPdesign:Dining Philosophers Problem}展示了
经典\IX{Dining Philosophers problem}的示意图~\cite{Dijkstra1971HOoSP}。
这个问题描述的是五位哲学家，他们除了思考和吃饭什么都不做，
而且吃的是一种"非常难对付的意大利面"，需要两把叉子才能
吃。\footnote{
	你也可以用筷子来类比思考这个问题。}
每位哲学家只能使用紧邻自己左右两侧的叉子，
而且在吃饱之前不会放下已经拿起的叉子。

\begin{figure*}
\centering
\IfTwoColumn{
\resizebox{5in}{!}{\includegraphics{cartoons/Dining-philosophers}}
}{
\resizebox{\onecolumntextwidth}{!}{\includegraphics{cartoons/Dining-philosophers}}
}
\caption{部分饥饿同样糟糕}
\ContributedBy{fig:cpu:Partial Starvation Is Also Bad}{Melissa Broussard}
\end{figure*}

目标是构造一个算法，从字面意义上防止\IX{starvation}（饥饿）。
一种饥饿场景是所有哲学家同时拿起了左手边的叉子。
因为在吃完之前没有人会放下自己的叉子，而且在至少有一人吃完之前
没有人能拿起第二把叉子，所以他们全都会饿死。
请注意，仅让至少一位哲学家能吃到东西是不够的。
如\cref{fig:cpu:Partial Starvation Is Also Bad}所示，
即使只有少数哲学家挨饿也是应当避免的。

\begin{figure}
\centering
\includegraphics[scale=.7]{SMPdesign/DiningPhilosopher5TB}
\caption{哲学家就餐问题，教科书解法}
\ContributedBy{fig:SMPdesign:Dining Philosophers Problem; Textbook Solution}{Kornilios Kourtis}
\end{figure}

\pplsur{Edsger W.}{Dijkstra}的解法使用了一个全局信号量，
在假设通信延迟可以忽略不计的情况下运作良好，
但这一假设在1980年代末或1990年代初就已失效。\footnote{
	站在2021年的角度来贬低Dijkstra，也就是事后50多年来看，
	是非常容易的。
	如果你仍然觉得有必要贬低Dijkstra，我的建议是先发表点什么，
	等50年，然后看看\emph{你的}想法能否经受住时间的考验。}
更近期的解法对叉子进行编号，如
\cref{fig:SMPdesign:Dining Philosophers Problem; Textbook Solution}所示。
每位哲学家先拿起自己盘子旁边编号较小的叉子，然后再拿另一把。
因此，图中最上方位置的哲学家先拿左手边的叉子，再拿右手边的叉子，
而其余哲学家则先拿右手边的叉子。
由于会有两位哲学家试图先拿起叉子~1，而且只有一位能成功，
所以四位哲学家将有五把叉子可用。
这四位中至少有一位将拥有两把叉子，从而能够进餐。

这种对资源编号并按编号顺序获取的通用技术被广泛用作
死锁预防技术。
然而，很容易设想出一种事件序列，使得即使所有人都饿了，
每次也只有一位哲学家在进餐：

\begin{enumerate}
    \item P2拿起叉子~1，阻止了P1拿到叉子。
    \item P3拿起叉子~2。
    \item P4拿起叉子~3。
    \item P5拿起叉子~4。
    \item P5拿起叉子~5并进餐。
    \item P5放下叉子~4和5。
    \item P4拿起叉子~4并进餐。
\end{enumerate}

简而言之，即使有足够的叉子让两位哲学家同时进餐，
这个算法也可能导致在所有五位哲学家都饿的情况下，
在任何给定时刻只有一位哲学家在进餐。
我们应该能做得更好！

\begin{figure}
\centering
\includegraphics[scale=.7]{SMPdesign/DiningPhilosopher4part-b}
\caption{哲学家就餐问题，分区方案}
\ContributedBy{fig:SMPdesign:Dining Philosophers Problem; Partitioned}{Kornilios Kourtis}
\end{figure}

一种方法如\cref{fig:SMPdesign:Dining Philosophers Problem; Partitioned}所示，
其中包含四位而非五位哲学家，以更好地说明分区技术。
这里上方和右方的哲学家共享一对叉子，
而下方和左方的哲学家共享另一对叉子。
如果所有哲学家同时饥饿，至少有两位总是能同时进餐。
此外，如图所示，叉子现在可以捆绑在一起，
一对同时拿起和放下，从而简化了获取和释放算法。

\QuickQuizSeries{%
\QuickQuizB{
	哲学家就餐问题有更好的解法吗？
}\QuickQuizAnswerB{
	一种改进的解法如
	\cref{fig:SMPdesign:Dining Philosophers Problem; Fully Partitioned}所示，
	即简单地为哲学家们额外提供五把叉子。
	现在五位哲学家可以同时进餐，
	而且永远不需要互相等待。
	此外，这种方法还大大改善了卫生状况。

\begin{figure}
\centering
\includegraphics[scale=.7]{SMPdesign/DiningPhilosopher5PEM}
\caption{哲学家就餐问题，完全分区}
\QContributedBy{fig:SMPdesign:Dining Philosophers Problem; Fully Partitioned}{Kornilios Kourtis}
\end{figure}

	这种解法在某些人看来可能像是作弊，但这种
	"作弊"正是找到许多并发问题优秀解法的关键，
	任何饥饿的哲学家都会同意这一点。

	这是\cpageref{sec:SMPdesign:Problems Dining Philosophers}
	提出的哲学家就餐并发进食问题的一个解法。
}\QuickQuizEndB
%
\QuickQuizE{
	你会如何验证一个声称能解决哲学家就餐问题的算法？
}\QuickQuizAnswerE{
	这很大程度上取决于算法的细节，但以下是几个着手点。

	首先，对于拿起左手和右手叉子是分开操作的算法，
	从所有叉子都在桌上开始。
	然后让所有哲学家尝试拿起第一把叉子。
	一旦所有哲学家要么拿到了第一把叉子，要么在等待别人放下第一把叉子，
	就让每位未等待的哲学家拿起第二把叉子。
	在任何无饥饿的解法中，此时至少有一位哲学家在进餐。
	如果有等待中的哲学家，重复此测试，
	最好在时序上加入随机变化。

	其次，创建一个压力测试，让哲学家在随机时间开始和停止进餐。
	生成饥饿和公平性条件并验证这些条件是否满足。
	以下是几个饥饿和公平性条件的示例：

	\begin{enumerate}
	\item	如果自某位哲学家尝试拿起某把叉子以来，
		其他所有哲学家都已停止进餐$N$次，
		则该哲学家应该已经成功拿起了那把叉子。
		对于使用高质量锁原语（或高质量原子操作）
		的高质量解法，$N=1$是可行的。
	\item	给定任何哲学家持有两把叉子到放下之间时间的上界$T$，
		任何哲学家的最大等待时间应该由$NT$限定，
		其中$N$不应比哲学家的数量大太多。
	\item	生成某种统计量，表示从哲学家尝试拿起第一把叉子
		到开始进餐之间的时间。
		该统计量越小，解法越好。
		均值、中位数和最大值都是有用的统计量，
		但检查完整的分布也很有启发性。
	\end{enumerate}

	鼓励读者实际测试本书中介绍的任何解法，
	特别是测试自己设计的解法。
}\QuickQuizEndE
}

这是"水平并行"~\cite{Inman85}
或"data parallelism"的一个例子，
之所以这样命名，是因为各对哲学家之间没有依赖关系。
在水平并行的数据处理系统中，给定的数据项
只会被一组复制的软件组件中的一个所处理。

\QuickQuiz{
	这种"水平并行"在什么意义上可以说是"水平"的？
}\QuickQuizAnswer{
	Jack Inman~\cite{Inman85}
	当时在处理协议栈，协议栈通常被垂直描绘，
	应用程序在顶部，硬件互连在底部。
	数据在这个栈中上下流动。
	"水平并行"并行处理来自不同网络连接的数据包，
	而"垂直并行"则并行处理给定数据包的不同协议处理步骤。

	"垂直并行"也称为"pipelining"。
}\QuickQuizEnd

\subsection{双端队列}
\label{sec:SMPdesign:Double-Ended Queue}

双端队列是一种数据结构，包含一个元素列表，
可以从任一端插入或移除元素~\cite{Knuth73}。
有人声称，允许在双端队列两端并发操作的基于锁的实现
是困难的~\cite{DanGrossman2007TMGCAnalogy}。
本节展示了分区设计策略如何产生
一个相当简单的实现，在以下小节中考察三种
通用方法。
但首先，我们应该如何验证一个并发双端队列？

\subsubsection{双端队列验证}
\label{sec:SMPdesign:Double-Ended Queue Validation}

一个好的起点是不变量。
例如，如果元素从双端队列的一端推入而从另一端弹出，
这些元素的顺序必须保持不变。
类似地，如果元素从队列的一端推入又从同一端弹出，
这些元素的顺序必须反转。
从队列中弹出的任何元素都必须是最近推入该队列的，
而且如果队列被清空，所有推入的元素都必须已经被弹出。

并发双端队列的测试套件（"\path{deqtorture.h}"）
起步阶段提供以下检查：

\begin{enumerate}
\item	由\co{CHECK_SEQUENCE_PAIR()}提供的元素顺序检查。
\item	由\co{melee()}提供的检查，验证弹出的元素是最近推入的。
\item	同样由\co{melee()}提供的检查，验证推入的元素在队列清空前已被弹出。
\end{enumerate}

该套件包含顺序测试和并发测试。
虽然该套件对于教科书代码来说是好且足够的，
但对于用于生产的代码，你应该进行更加彻底的测试。
\Cref{chp:Validation,chp:Formal Verification}涵盖了大量的
验证工具和技术。

有了原型测试套件，我们就可以在下面的小节中
考察双端队列算法了。

\subsubsection{左右手锁}
\label{sec:SMPdesign:Left- and Right-Hand Locks}

\begin{figure}
\centering
\resizebox{3in}{!}{\includegraphics{SMPdesign/lockdeq}}
\caption{带左右手锁的双端队列}
\label{fig:SMPdesign:Double-Ended Queue With Left- and Right-Hand Locks}
\end{figure}

一种看似直接的方法是使用双向链表，
配备一个左手锁用于左端的入队和出队操作，
以及一个右手锁用于右端操作，
如\cref{fig:SMPdesign:Double-Ended Queue With Left- and Right-Hand Locks}所示。
然而，这种方法的问题在于，当列表中少于四个元素时，
两个锁的域必须重叠。
这种重叠是因为移除任何给定元素不仅影响该元素本身，
还影响其左右邻居。
这些域在图中用颜色表示，蓝色带向下条纹表示
左手锁的域，红色带向上条纹表示右手锁的域，
紫色（无条纹）表示重叠域。
虽然可以创建一个以这种方式工作的算法，
也许使用类似于Michael和Scott提出的双锁队列的虚拟元素~\cite{MichaelScott96}，
但它至少有五种特殊情况这一事实应该引起高度警惕，
特别是考虑到列表另一端的并发活动可以随时
将队列从一种特殊情况转换到另一种。

最好考虑其他设计。

\subsubsection{复合双端队列}
\label{sec:SMPdesign:Compound Double-Ended Queue}

\begin{figure}
\centering
\resizebox{3in}{!}{\includegraphics{SMPdesign/lockdeqpair}}
\caption{复合双端队列}
\label{fig:SMPdesign:Compound Double-Ended Queue}
\end{figure}

\cref{fig:SMPdesign:Compound Double-Ended Queue}展示了
一种强制锁域不重叠的方法。
两个独立的双端队列串联运行，每个都由自己的锁保护。
这意味着元素有时必须在两个双端队列之间转移，
这种情况下必须同时持有两个锁。
可以使用简单的锁层次来避免\IX{deadlock}，例如，
始终在获取右手锁之前先获取左手锁。
这将比对同一个双端队列应用两个锁简单得多，
因为我们可以无条件地将元素左入队到左手队列，
将元素右入队到右手队列。
主要的复杂性出现在从空队列出队时，
这种情况下需要：

\begin{enumerate}
\item	如果持有右手锁，释放它并获取左手锁。
\item	获取右手锁。
\item	在两个队列之间重新平衡元素。
\item	如果有元素，则移除所需的元素。
\item	释放两个锁。
\end{enumerate}

\QuickQuiz{
	在这个复合双端队列实现中，如果在释放
	并重新获取锁的过程中队列变为非空，应该怎么办？
}\QuickQuizAnswer{
	在这种情况下，只需从非空队列出队一个元素，
	释放两个锁，然后返回。
}\QuickQuizEnd

生成的代码（\path{locktdeq.c}）相当直观。
重新平衡操作可能会使给定元素在两个队列之间来回转移，
浪费时间，可能需要依赖工作负载的启发式方法来获得最佳性能。
虽然这在某些情况下可能是最好的方法，
但尝试一种具有更高确定性的算法也很有趣。

\subsubsection{哈希双端队列}
\label{sec:SMPdesign:Hashed Double-Ended Queue}

确定性地分区数据结构的最简单和最有效的方法之一是对其进行哈希。
可以通过根据元素在列表中的位置为其分配序列号来简单地对双端队列进行哈希，
使得第一个左入队到空队列的元素编号为零，
第一个右入队到空队列的元素编号为一。
一系列左入队到原本空闲队列的元素将被分配递减的编号（$-1$, $-2$, $-3$, \ldots），
而一系列右入队到原本空闲队列的元素将被分配递增的编号（2, 3, 4, \ldots）。
关键点是实际上不需要表示给定元素的编号，
因为这个编号将由其在队列中的位置隐含。

\begin{figure}
\centering
\resizebox{3in}{!}{\includegraphics{SMPdesign/lockdeqhash}}
\caption{哈希双端队列}
\label{fig:SMPdesign:Hashed Double-Ended Queue}
\end{figure}

采用这种方法，我们分配一个锁来保护左手索引，
一个锁来保护右手索引，以及每个哈希链一个锁。
\Cref{fig:SMPdesign:Hashed Double-Ended Queue}展示了
给定四个哈希链的结果数据结构。
注意锁域不重叠，并且通过在链锁之前获取索引锁，
以及永远不同时获取多个给定类型（索引或链）的锁，
来避免死锁。

\begin{figure}
\centering
\resizebox{3in}{!}{\includegraphics{SMPdesign/lockdeqhash1R}}
\caption{插入后的哈希双端队列}
\label{fig:SMPdesign:Hashed Double-Ended Queue After Insertions}
\end{figure}

每个哈希链本身就是一个双端队列，在这个例子中，
每个持有每四个元素。
\Cref{fig:SMPdesign:Hashed Double-Ended Queue After Insertions}
的最上部分显示了单个元素（"R$_1$"）被右入队后的状态，
右手索引已递增以引用哈希链~2。
同一图的中间部分显示了再右入队三个元素后的状态。
如你所见，索引回到了初始状态
（见\cref{fig:SMPdesign:Hashed Double-Ended Queue}），
然而，每个哈希链现在都非空。
该图的下部分显示了三个额外元素被左入队
以及一个额外元素被右入队后的状态。

从\cref{fig:SMPdesign:Hashed Double-Ended Queue After Insertions}
显示的最后状态开始，
左出队操作将返回元素"L$_{-2}$"并使
左手索引引用哈希链~2，该链随后
将只包含一个元素（"R$_2$"）。
在这种状态下，左入队与右入队并发运行
将导致\IX{lock contention}，但可以通过使用更大的哈希表
将这种竞争的概率降低到任意低的水平。

\begin{figure}
\centering
\resizebox{1.5in}{!}{\includegraphics{SMPdesign/lockdeqhashlots}}
\caption{包含16个元素的哈希双端队列}
\label{fig:SMPdesign:Hashed Double-Ended Queue With 16 Elements}
\end{figure}

\Cref{fig:SMPdesign:Hashed Double-Ended Queue With 16 Elements}
展示了16个元素如何在四个哈希桶的
并行双端队列中组织。
每个底层的单锁双端队列持有完整并行双端队列的
四分之一切片。

\begin{listing}
\input{CodeSamples/SMPdesign/lockhdeq@struct_pdeq.fcv}
\caption{基于锁的并行双端队列数据结构}
\label{lst:SMPdesign:Lock-Based Parallel Double-Ended Queue Data Structure}
\end{listing}

\Cref{lst:SMPdesign:Lock-Based Parallel Double-Ended Queue Data Structure}
展示了对应的C语言数据结构，假设已存在一个提供简单加锁
双端队列实现的\co{struct deq}。
\begin{fcvref}[ln:SMPdesign:lockhdeq:struct_pdeq]
该数据结构包含\clnref{llock}上的左手锁，
\clnref{lidx}上的左手索引，\clnref{rlock}上的右手锁
（在实际实现中是cache对齐的），
\clnref{ridx}上的右手索引，以及最后是\clnref{bkt}上的
简单基于锁的双端队列哈希数组。
高性能实现当然会使用填充或特殊的
对齐指令来避免\IX{false sharing}。
\end{fcvref}

\begin{listing}
\input{CodeSamples/SMPdesign/lockhdeq@pop_push.fcv}
\caption{基于锁的并行双端队列实现}
\label{lst:SMPdesign:Lock-Based Parallel Double-Ended Queue Implementation}
\end{listing}

\Cref{lst:SMPdesign:Lock-Based Parallel Double-Ended Queue Implementation}
（\path{lockhdeq.c}）
展示了入队和出队函数的实现。\footnote{
	人们可以轻松地用任意数量的语言创建多态实现，
	但这留作读者练习。}
讨论将集中于左手操作，因为右手操作
可以简单地从左手操作推导出来。

\begin{fcvref}[ln:SMPdesign:lockhdeq:pop_push:popl]
\Clnrefrange{b}{e}展示了\co{pdeq_pop_l()}，
它左出队并返回一个元素（如果可能的话），否则返回\co{NULL}。
\Clnref{acq}获取左手自旋锁，
\clnref{idx}计算要出队的索引。
\Clnref{deque}出队元素，如果\clnref{check}发现结果
非\co{NULL}，\clnref{record}记录新的左手索引。
无论哪种情况，\clnref{rel}释放锁，最后\clnref{return}
返回元素（如果有的话），否则返回\co{NULL}。
\end{fcvref}

\begin{fcvref}[ln:SMPdesign:lockhdeq:pop_push:pushl]
\Clnrefrange{b}{e}展示了\co{pdeq_push_l()}，
它左入队指定的元素。
\Clnref{acq}获取左手锁，
\clnref{idx}取得左手索引。
\Clnref{enque}将指定元素左入队到
由左手索引索引的双端队列上。
\Clnref{update}然后更新左手索引，
\clnref{rel}释放锁。
\end{fcvref}

如前所述，右手操作与其左手对应操作完全类似，
因此其分析留作读者练习。

\QuickQuiz{
	哈希双端队列是一个好的解决方案吗？
	为什么是或者为什么不是？
}\QuickQuizAnswer{
	回答这个问题的最好方法是在多种不同的多处理器系统上
	运行\path{lockhdeq.c}，强烈鼓励你这样做。
	一个令人担忧的原因是，此实现上的每个操作
	必须获取的不是一个而是两个锁。
	% Getting about 500 nanoseconds per element when used as
	% a queue on a 4.2GHz Power system.  This is roughly the same as
	% the version covered by a single lock.  Sequential (unlocked
	% variant is more than an order of magnitude faster!

	第一个设计良好的性能研究将被引用。\footnote{
		Dalessandro等人的研究~\cite{LukeDalessandro:2011:ASPLOS:HybridNOrecSTM:deque}
		和Dice等人的研究~\cite{DavidDice:2010:SCA:HTM:deque}是
		很好的起点。}
	不要忘记与顺序实现进行比较！
}\QuickQuizEnd

\subsubsection{复合双端队列再探}
\label{sec:SMPdesign:Compound Double-Ended Queue Revisited}

本节重新审视复合双端队列，使用一种简单的重新平衡方案，
将所有元素从非空队列移动到变空的队列。

\QuickQuiz{
	将\emph{所有}元素移动到变空的队列？
	在什么可能的世界里，这种愚蠢的解决方案在任何意义上是最优的？？？
}\QuickQuizAnswer{
	当数据流方向很少切换时，它是最优的。
	当然，如果双端队列同时从两端被清空，
	这将是一个极其糟糕的选择。
	这自然引出了另一个问题，即在什么可能的世界里
	同时从两端清空是合理的做法。
	工作窃取队列是这个问题的一个可能答案。
}\QuickQuizEnd

与上一节介绍的哈希实现不同，复合实现将建立在
一个既不使用锁也不使用原子操作的双端队列顺序实现之上。

\begin{listing}
\ebresizeverb{.97}{\input{CodeSamples/SMPdesign/locktdeq@pop_push.fcv}}
\caption{复合并行双端队列实现}
\label{lst:SMPdesign:Compound Parallel Double-Ended Queue Implementation}
\end{listing}

\Cref{lst:SMPdesign:Compound Parallel Double-Ended Queue Implementation}
展示了该实现。
与哈希实现不同，这个复合实现是不对称的，
因此我们必须分别考虑\co{pdeq_pop_l()}
和\co{pdeq_pop_r()}的实现。

\QuickQuiz{
	为什么复合并行双端队列的实现不能是对称的？
}\QuickQuizAnswer{
	需要通过施加锁层次来避免死锁，
	这迫使了不对称性，就像哲学家就餐问题中叉子编号解法
	的情况一样
	（见\cref{sec:SMPdesign:Dining Philosophers Problem}）。
}\QuickQuizEnd

\begin{fcvref}[ln:SMPdesign:locktdeq:pop_push:popl]
\co{pdeq_pop_l()}的实现如图的
\clnrefrange{b}{e}所示。
\Clnref{acq:l}获取左手锁，
\clnref{rel:l}释放它。
\Clnref{deq:ll}尝试从左手底层双端队列左出队一个元素，
如果成功，则跳过\clnrefrange{acq:r}{skip}直接返回此元素。
否则，\clnref{acq:r}获取右手锁，\clnref{deq:lr}
从右手队列左出队一个元素，
\clnref{move}将右手队列上剩余的元素
移动到左手队列，\clnref{init:r}初始化右手队列，
\clnref{rel:r}释放右手锁。
在\clnref{deq:lr}上出队的元素（如果有的话）将被返回。
\end{fcvref}

\begin{fcvref}[ln:SMPdesign:locktdeq:pop_push:popr]
\co{pdeq_pop_r()}的实现如图的\clnrefrange{b}{e}所示。
与之前一样，\clnref{acq:r1}获取右手锁
（\clnref{rel:r2}释放它），\clnref{deq:rr1}尝试从右手队列
右出队一个元素，如果成功，
则跳过\clnrefrange{rel:r1}{skip2}直接返回此元素。
然而，如果\clnref{check1}确定没有可出队的元素，
\clnref{rel:r1}释放右手锁，
\clnrefrange{acq:l}{acq:r2}按正确顺序获取两个锁。
\Clnref{deq:rr2}然后再次尝试从右手列表右出队一个元素，
如果\clnref{check2}确定第二次尝试也失败了，
\clnref{deq:rl}从左手队列右出队一个元素
（如果有可用的），\clnref{move}将左手队列的
剩余元素移动到右手队列，\clnref{init:l}
初始化左手队列。
无论哪种情况，\clnref{rel:l}释放左手锁。
\end{fcvref}

\QuickQuizSeries{%
\QuickQuizB{
	为什么有必要在
	\cref{lst:SMPdesign:Compound Parallel Double-Ended Queue Implementation}的
	\clnrefr{ln:SMPdesign:locktdeq:pop_push:popr:deq:rr2}
	上重试右出队操作？
}\QuickQuizAnswerB{
	\begin{fcvref}[ln:SMPdesign:locktdeq:pop_push:popr]
	这个重试是必要的，因为在本线程在\clnref{rel:r1}上放下
	\co{d->rlock}和在\clnref{acq:r2}上重新获取同一个锁之间，
	其他某个线程可能已经入队了一个元素。
	\end{fcvref}
}\QuickQuizEndB
%
\QuickQuizE{
	左手锁肯定\emph{有时候}是可用的！！！
	那么为什么
	\cref{lst:SMPdesign:Compound Parallel Double-Ended Queue Implementation}的
	\clnrefr{ln:SMPdesign:locktdeq:pop_push:popr:rel:r1}
	必须无条件释放右手锁？
}\QuickQuizAnswerE{
	可以使用\co{spin_trylock()}在左手锁可用时
	尝试获取它。
	然而，失败的情况仍然需要放下右手锁，
	然后按顺序重新获取两个锁。
	进行这个转换（以及确定是否值得）
	留作读者练习。
}\QuickQuizEndE
}

\begin{fcvref}[ln:SMPdesign:locktdeq:pop_push:pushl]
\co{pdeq_push_l()}的实现如
\cref{lst:SMPdesign:Compound Parallel Double-Ended Queue Implementation}的
\clnrefrange{b}{e}所示。
\Clnref{acq:l}获取左手自旋锁，
\clnref{que:l}将元素左入队到左手队列，
最后\clnref{rel:l}释放锁。
\end{fcvref}
\begin{fcvref}[ln:SMPdesign:locktdeq:pop_push:pushr]
\co{pdeq_push_r()}的实现（如\clnrefrange{b}{e}所示）
非常类似。
\end{fcvref}

\QuickQuiz{
	但在数据仅沿一个方向流动的情况下，
	\cref{lst:SMPdesign:Compound Parallel Double-Ended Queue Implementation}
	中展示的算法会使两端在消费端清空其底层
	双端队列时都尝试获取同一个锁。
	这难道不意味着有时即使队列包含任意大量的元素，
	该算法也无法提供对队列两端的并发访问吗？
}\QuickQuizAnswer{
	确实如此！

	但其他声称具有这一属性的算法也是如此。
	例如，在使用基于哈希锁数组的软件事务内存
	机制的解法中，最左端和最右端元素的地址
	有时恰好会哈希到同一个锁。
	这些哈希冲突也会阻止并发访问。
	再举一个例子，使用带有软件回退的硬件事务内存
	机制的解法~\cite{Yoo:2013:PEI:2503210.2503232,RickMerrit2011PowerTM,ChristianJacobi2012MainframeTM}
	通常在这些软件回退中使用锁，因此
	会受到（虽然希望很少）这些锁方案所具有的
	任何并发限制的影响。

	因此，截至2021年，所有并发双端队列问题的实际解法
	至少在某些情况下都无法提供完全的并发性，
	包括复合双端队列。
}\QuickQuizEnd

\subsubsection{双端队列讨论}
\label{sec:SMPdesign:Double-Ended Queue Discussion}

复合实现比\cref{sec:SMPdesign:Hashed Double-Ended Queue}中介绍的
哈希变体稍微复杂一些，但仍然相当简单。
当然，更智能的重新平衡方案可以任意复杂，
但这里展示的简单方案已被证明与软件替代方案相比
性能良好~\cite{LukeDalessandro:2011:ASPLOS:HybridNOrecSTM:deque}，
甚至与使用硬件辅助的算法相比也是
如此~\cite{DavidDice:2010:SCA:HTM:deque}。
然而，从这种方案中我们最多能期望的是2倍的可扩展性，
因为最多只有两个线程可以同时持有出队的锁。
这一限制同样适用于基于\IXacrl{nbs}的算法，
例如Michael的基于compare-and-swap的出队
算法~\cite{DBLP:conf/europar/Michael03}。\footnote{
	这篇论文很有趣，因为它表明不需要特殊的
	double-compare-and-swap（DCAS）指令来实现双端队列的无锁实现。
	相反，通用的compare-and-swap（例如x86 cmpxchg）
	就足够了。}

\QuickQuiz{
	为什么双端队列问题不是有一个而是有两个解法？
}\QuickQuizAnswer{
	实际上至少有三个。
	第三个由Dominik Dingel提出，有趣地使用了
	读写锁，可以在\path{lockrwdeq.c}中找到。

	因此在\cpageref{sec:SMPdesign:Problems Double-Ended Queue}上，
	基于锁的双端队列问题不是有一个，而是有三个解法！
}\QuickQuizEnd

事实上，正如Dice等人~\cite{DavidDice:2010:SCA:HTM:deque}所指出的，
一个未同步的单线程双端队列在性能上显著
优于他们研究的任何并行实现。
因此，关键点是无论实现方式如何，
向共享队列入队或出队可能会有显著的开销。
鉴于\cref{chp:Hardware and its Habits}中的材料
以及这些队列严格的先进先出（FIFO）特性，
这应该不足为奇。

此外，这些严格的FIFO队列仅在
\emph{线性化点}~\cite{Herlihy:1990:LCC:78969.78972}\footnote{
	简而言之，线性化点是给定函数中的一个单一点，
	在该点上可以认为该函数已经生效。
	在这个基于锁的实现中，线性化点
	可以说是在完成工作的临界区内的任何位置。}
方面是严格FIFO的，而这些线性化点对调用者不可见，
实际上，在这些例子中，线性化点被埋在基于锁的
临界区中。
这些队列在（比如说）各个操作开始的时间
方面不是严格FIFO的~\cite{AndreasHaas2012FIFOisnt}。
这表明严格的FIFO属性在并发程序中并不那么有价值，
事实上，Kirsch等人提出了限制较少的队列，
提供了改进的性能和
可扩展性~\cite{ChristophMKirsch2012FIFOisntTR}。\footnote{
	Nir Shavit出于大致相同的原因
	提出了宽松的栈~\cite{Shavit:2011:DSM:1897852.1897873}。
	这种情况使一些人认为线性化点对理论家比对开发者更有用，
	也使另一些人怀疑这些数据结构和算法的设计者
	在多大程度上考虑了用户的需求。}
总之，如果你将并发程序使用的所有数据都通过
一个单一队列推送，你真的需要重新思考你的整体设计。

\subsection{分区示例讨论}
\label{sec:SMPdesign:Partitioning Example Discussion}

在\cref{sec:SMPdesign:Dining Philosophers Problem}中
Quick Quiz的答案中给出的哲学家就餐问题最优解
是"水平并行"或"data parallelism"的一个极好的例子。
在这种情况下，同步开销几乎（甚至完全）为零。
相比之下，双端队列实现是"垂直并行"或
"pipelining"的例子，因为数据从一个线程移动到另一个线程。
流水线所需的更紧密协调反过来要求
更大的工作单元来获得给定水平的\IX{efficiency}。

\QuickQuizSeries{%
\QuickQuizB{
	串联双端队列的速度大约是哈希双端队列的两倍，
	即使我将哈希表的大小增加到一个疯狂的巨大数字。
	这是为什么？
}\QuickQuizAnswerB{
	哈希双端队列的锁设计只允许每端一次一个线程，
	而且每个操作需要两次锁获取。
	串联双端队列同样允许每端一次一个线程，
	但在常见情况下每个操作只需要一次锁获取。
	因此，串联双端队列应该预期会
	优于哈希双端队列。

	你能创建一个允许每端多个并发操作的双端队列吗？
	如果能，怎么做？
	如果不能，为什么不能？
}\QuickQuizEndB
%
\QuickQuizE{
	有没有明显更好的方法来处理双端队列的并发性？
}\QuickQuizAnswerE{
	一种方法是转换要解决的问题，
	使得可以并行使用多个双端队列，
	允许使用更简单的单锁双端队列，
	并且也许还可以用一对传统的单端队列
	替换每个双端队列。
	如果没有这种"水平扩展"，加速比限制在2.0。
	相比之下，水平扩展设计可以实现非常大的加速比，
	如果有多个线程在队列的任一端工作则特别有吸引力，
	因为在这种多线程情况下，出队
	根本无法提供强排序保证。
	毕竟，一个给定线程先移除了一个元素
	并不意味着它会先处理该
	元素~\cite{AndreasHaas2012FIFOisnt}。
	如果没有保证，我们不如获取
	拒绝提供这些保证所带来的性能优势。
	% about twice as fast as hashed version on 4.2GHz Power.

	无论问题是否可以转换为使用多个队列，
	都值得问一下是否可以批处理工作，
	使每个入队和出队操作对应更大的工作单元。
	这种批处理方法减少了对队列数据结构的竞争，
	从而提高了性能和可扩展性，
	如\cref{sec:SMPdesign:Synchronization Granularity}中
	所示。
	毕竟，如果你必须承受高同步开销，
	请确保你获得了相应的回报。

	其他研究人员正在研究利用队列中
	有限排序保证的其他方法~\cite{ChristophMKirsch2012FIFOisntTR}。
}\QuickQuizEndE
}

这两个例子展示了分区在设计并行算法方面有多么强大。
\Cref{sec:SMPdesign:Locking Granularity and Performance}
简要考察了第三个例子——矩阵乘法。
然而，这三个例子都迫切需要更多更好的并行程序设计准则，
这是下一节要探讨的主题。
